\documentclass[11pt,a4paper]{article}
\pdfoutput=1
\usepackage[T1]{fontenc}
\usepackage{lmodern,amsmath,amssymb,amsthm}
\usepackage[margin=28mm]{geometry}
\usepackage{microtype}
\usepackage[hidelinks]{hyperref}
\newcommand{\Prob}{\mathbb P}
\newcommand{\E}{\mathbb E}
\newcommand{\dd}{\,\mathrm d}
\newtheorem{theorem}{Theorem}
\title{Extreme excursions and terminal-sign retention in stationary Brownian resetting}
\author{Xiaojian Zeng\\
\small PBC School of Finance, Tsinghua University, 100190\\
\small \texttt{xiaojianzeng@foxmail.com}}
\date{}
\begin{document}
\maketitle
\begin{abstract}
For stationary Brownian motion with independent exponential resetting, we determine
the fixed-horizon probability of a high excursion followed by a positive terminal
position. The exponential tail coefficient contains the probability of retaining
a positive endpoint after the first reset, averaged over the remaining duration.
An explicit Gaussian-tail remainder controls the finite-threshold approximation
for every fixed positive horizon and reset rate and every finite reset position.
This establishes the conjectured joint-tail asymptotic of Dębicki, Hashorva and
Michna. When the reset position is zero, the coefficient is exactly midway between
their conjectured bounds. The argument separates the stationary initial excursion
from later reset-origin excursions and does not interchange terminal-threshold
and high-excursion limits.
\end{abstract}

\section{Observable and model}
Endpoint monitoring of a resetting diffusion need not retain evidence of an earlier
large excursion: a reset erases its position. We calculate this effect for a particle
in a nonequilibrium stationary state on the full real line. Dębicki, Hashorva and
Michna~\cite{DHM} obtained the stationary maximum tail and joint tails for terminal
thresholds proportional to the excursion level. Their Remark~3 leaves the zero
terminal threshold conjectural. We determine its coefficient and give an explicit
finite-threshold error. Renewal and two-time methods for complete and incomplete
resetting are developed in~\cite{TMM}; the contribution here is the stationary
path-maximum/terminal-sign observable, rather than a new renewal formalism.

Let $B$ be a standard Brownian motion, $\operatorname{Var}(B_t)=t$, and let $N$
be an independent Poisson process of rate $r>0$. Complete resets place the particle
at $a\in\mathbb R$:
\begin{equation}\label{eq:model}
Y_t=Y_0+B_t+\int_{(0,t]}(a-Y_{s-})\dd N_s,\qquad
p_\infty(y)=\frac{\alpha}{2}e^{-\alpha|y-a|},\quad \alpha=\sqrt{2r}.
\end{equation}
The initial position has density $p_\infty$ and is independent of future noise.
The diffusion coefficient is $D=1/2$ in these units; there is no drift or reflecting
boundary. The generator is
$\mathcal L f(y)=f''(y)/2+r[f(a)-f(y)]$.
The displayed density solves
$p_\infty''/2-rp_\infty+r\delta_a=0$ distributionally and has unit mass.
Lengths are expressed in units for which the Brownian diffusion length is
$\sqrt t$; thus $\alpha a$, $\alpha u$ and $\alpha\sqrt t$ are dimensionless.
Put $M_T=\sup_{0\leq t\leq T}Y_t$, with $T>0$ fixed.

For a process started deterministically at $a$, its positive-terminal probability
after duration $v>0$ is
\begin{equation}\label{eq:q}
q_a(v)=e^{-rv}\Phi(a/\sqrt v)
       +\int_0^v r e^{-rw}\Phi(a/\sqrt w)\dd w.
\end{equation}
Here $\Phi$ is the standard normal distribution function. The last-reset age has
an atom $e^{-rv}$ at $v$ and density $re^{-rw}$ on $(0,v)$, proving
\eqref{eq:q}. The strict event at zero duration gives $q_a(0)=\mathbf1_{\{a>0\}}$.
Its right limit is $1$, $1/2$ or $0$ for positive, zero or negative $a$,
respectively. Isolated endpoint values do not affect the integrals below.

\section{Coefficient and finite-threshold bound}
\begin{theorem}\label{thm:main}
For every fixed $T>0$, $r>0$ and finite $a\in\mathbb R$ in \eqref{eq:model},
\begin{equation}\label{eq:C}
\lim_{u\to\infty}e^{\alpha u}\Prob(M_T>u,Y_T>0)
=C(T,r,a)
=e^{\alpha a}\left[
\Phi(\alpha\sqrt T)+r\int_0^T\Phi(\alpha\sqrt s)q_a(T-s)\dd s\right].
\end{equation}
For $u>\max(a,0)$, let $d=u-a$. Then
\begin{equation}\label{eq:error}
\left|e^{\alpha u}\Prob(M_T>u,Y_T>0)-C(T,r,a)\right|
\leq R(u;T,r,a),
\end{equation}
where
\begin{align}\label{eq:R}
R(u;T,r,a)
={}&e^{\alpha a+rT}\Phi\!\left(\frac{\alpha T-d}{\sqrt T}\right)\nonumber\\
&+2e^{\alpha u}\left[
e^{-rT}\Phi\!\left(-\frac u{\sqrt T}\right)
+rT\Phi\!\left(-\frac d{\sqrt T}\right)\right].
\end{align}
The bound tends to zero for each fixed parameter triple.
\end{theorem}
In particular, $C$ is finite and positive and obeys the strict inequalities
\begin{equation}\label{eq:bounds}
e^{\alpha a}\Phi(\alpha\sqrt T)
<C(T,r,a)<
e^{\alpha a}\left[\Phi(\alpha\sqrt T)
+r\int_0^T\Phi(\alpha\sqrt s)\dd s\right].
\end{equation}
These establish the bounds conjectured in Remark~3 of~\cite{DHM}. No uniformity
in simultaneous limits of $T,r,a$ and $u$ is asserted.

\section{Proof}
\subsection{Separate the initial segment}
Let $\tau$ be the first future reset and let $E_u$ be a crossing of $u$
in the initial Brownian segment, stopped at $\tau$ or $T$.
Write $m_t=\sup_{0\leq s\leq t}B_s$ and
\[
H(t,u)=\Prob(Y_0+m_t>u),\qquad
j(t,u)=\Prob(Y_0+m_t>u,Y_0+B_t>0).
\]
Conditional on $\tau=s$, complete resetting makes the future independent of
the initial crossing, with terminal sign probability $q_a(T-s)$. Therefore
\begin{equation}\label{eq:F}
F(u):=\Prob(E_u,Y_T>0)
=e^{-rT}j(T,u)+\int_0^T re^{-rs}H(s,u)q_a(T-s)\dd s.
\end{equation}
Crossing at a reset instant has probability zero; continuity within the segment
also represents its pre-reset supremum. Since $u>a$, a reset jump cannot create
a crossing from below. Every crossing not counted in $E_u$ must occur in one
of the $N_T$ later Brownian segments starting at $a$.
Conditional on all reset times, each has length at most $T$. Reflection and
a union bound give
\begin{equation}\label{eq:later}
0\leq \Prob(M_T>u,Y_T>0)-F(u)
\leq 2rT\Phi(-d/\sqrt T).
\end{equation}
Indeed the conditional sum is at most $2N_T\Phi(-d/\sqrt T)$ and
$\E N_T=rT$. Independence between segment events is unnecessary. This treats
all reset counts without an infinite-series truncation.

\subsection{Control loss at the Brownian endpoint}
On $\{Y_0+m_T>u,\ Y_0+B_T\leq0\}$ one has $m_T-B_T>u$.
The process $B_{T-h}-B_T$, $0\leq h\leq T$, is a Brownian motion in distribution,
so $m_T-B_T$ has the distribution of a Brownian maximum. For $u>0$,
\begin{equation}\label{eq:loss}
0\leq H(T,u)-j(T,u)\leq2\Phi(-u/\sqrt T).
\end{equation}
The implication includes initial positions above $u$ and does not impose a sign
condition on the reset position.

\subsection{Control the stationary initial tail}
Let $X=Y_0-a$ and $L(z)=\Prob(X>z)$. Its centered Laplace survival is
\[
L(z)=
\begin{cases}
\tfrac12e^{-\alpha z},&z\geq0,\\
1-\tfrac12e^{\alpha z},&z<0.
\end{cases}
\]
For $d>0$ and $m\geq0$,
\begin{equation}\label{eq:laplace}
0\leq \frac12e^{\alpha m}-e^{\alpha d}L(d-m)
\leq\frac12e^{\alpha m}\mathbf1_{\{m>d\}}.
\end{equation}
The difference is zero for $m\leq d$; for $m>d$ its nonnegativity follows
from $\cosh(\alpha(m-d))\geq1$, and its upper bound from $L\geq0$.
Since $H(t,u)=\E L(d-m_t)$ and $m_t\leq m_T$ under the same Brownian path,
\begin{equation}\label{eq:Hbound}
0\leq A(t)-e^{\alpha d}H(t,u)\leq D_T(d),\qquad 0\leq t\leq T,
\end{equation}
where half-normal integration yields
\begin{align}\label{eq:AD}
A(t)&=\tfrac12\E e^{\alpha m_t}=e^{rt}\Phi(\alpha\sqrt t),\\
D_T(d)&=\tfrac12\E[e^{\alpha m_T}\mathbf1_{\{m_T>d\}}]
=e^{rT}\Phi((\alpha T-d)/\sqrt T).\nonumber
\end{align}
At $t=0$, $A(0)=1/2$ and $H(0,u)=e^{-\alpha d}/2$,
so the time endpoint is covered.

\subsection{Assemble the bounds}
In \eqref{eq:F}, replace $j(T,u)$ by $H(T,u)$ and then
$e^{\alpha u}H(t,u)$ by $e^{\alpha a}A(t)$.
Equation~\eqref{eq:AD} gives exactly \eqref{eq:C}. The initial-tail error is
at most $e^{\alpha a}D_T(d)$ times
\[
e^{-rT}+\int_0^T re^{-rs}q_a(T-s)\dd s\leq1.
\]
The endpoint-loss error from \eqref{eq:loss} is at most
$2e^{\alpha u-rT}\Phi(-u/\sqrt T)$.
The later-segment contribution in \eqref{eq:later} adds
$2rTe^{\alpha u}\Phi(-d/\sqrt T)$. Their sum is \eqref{eq:R}.
Each Gaussian tail vanishes after multiplication by a fixed exponential
in $u$, proving \eqref{eq:error} and the limit. Finally $0<q_a(v)<1$
for $v>0$ and finite $a$, proving \eqref{eq:bounds}.
\hfill$\square$

\section{Evaluation and physical interpretation}
\subsection{Controlled evaluation}
Formula~\eqref{eq:C} involves only finite one-dimensional integrals.
For exact-real enclosures, define $f_a(w)=\Phi(a/\sqrt w)$ using its right
limit at $w=0$. On a partition $0=w_0<\cdots<w_n=v$, each contribution
to the integral in \eqref{eq:q} is between
\[
(e^{-rw_i}-e^{-rw_{i+1}})
\min\{f_a(w_i),f_a(w_{i+1})\}
\quad\hbox{and the corresponding maximum.}
\]
Monotonicity of $f_a$ gives an enclosure width no greater than
$r\max_i(w_{i+1}-w_i)$. Add the atom $e^{-rv}f_a(v)$.
Differentiating \eqref{eq:q} gives $q_a'(v)=e^{-rv}f_a'(v)$, so $q_a$
is decreasing for $a>0$, increasing for $a<0$, and constant for $a=0$.
For the outer integral, endpoint bounds on both nonnegative factors bracket
their product on every subinterval. Use the right limit of $q_a$ at zero
for these endpoint bounds. Refinement converges because the factors have bounded
variation and continuous extensions from positive durations.
Certified evaluations of $\Phi$ and the exponential, or directed rounding,
are required to turn these exact-real bounds into machine certificates.

An alternative evaluation eliminates the inner quadrature. For $b=|a|>0$, put
\begin{equation}\label{eq:G}
G(v)=e^{-\alpha b}\Phi(\alpha\sqrt v-b/\sqrt v)
+e^{\alpha b}\Phi(-\alpha\sqrt v-b/\sqrt v).
\end{equation}
Its derivative is
$G'(v)=e^{-rv}b e^{-b^2/(2v)}/\sqrt{2\pi v^3}$ and $G(0+)=0$.
Consequently $q_a(v)=1-G(v)/2$ for $a>0$ and $q_a(v)=G(v)/2$
for $a<0$, by the derivative and initial limit of \eqref{eq:q}.
Large exponential factors in \eqref{eq:G} require stable special-function
evaluation; the bounded-integral construction avoids that difficulty.

\subsection{Zero reset position and retention}
When $a=0$, symmetry gives $q_0(v)=1/2$ for $v>0$:
\begin{equation}\label{eq:zero}
C(T,r,0)=\Phi(\alpha\sqrt T)+\frac r2
\int_0^T\Phi(\alpha\sqrt s)\dd s.
\end{equation}
This is the midpoint of the two bounding coefficients. With $k=\alpha\sqrt T$
and $\phi=\Phi'$,
\begin{equation}\label{eq:I}
\int_0^T\Phi(\alpha\sqrt s)\dd s
=T\Phi(k)+\frac{k\phi(k)-\Phi(k)+1/2}{\alpha^2}.
\end{equation}
For $T=r=1$, \eqref{eq:zero} gives $C\simeq1.32857643$.
This decimal is an ordinary floating-point evaluation of an exact expression.

The known unconditional coefficient in Theorem~5 of~\cite{DHM} is the upper
expression in \eqref{eq:bounds}. Dividing $C$ by that coefficient gives the
limiting fraction of extreme paths retained by the positive endpoint monitor.
The initial high segment controls the exponential scale; later excursions
starting at a finite reset point have Gaussian-small crossing probabilities.
After an initial excursion, a subsequent reset can erase the extreme position,
but its remaining evolution retains the endpoint sign with probability $q_a$.
The first term in \eqref{eq:C} is the no-future-reset contribution; the integral
quantifies this additional post-reset retention.

Changing $a$ alters both the stationary tail multiplier and the post-reset sign
probability, because the monitoring boundary stays at zero. Thus translation
cannot be represented by the multiplier $e^{\alpha a}$ alone.
The fixed zero threshold is evaluated directly: taking a terminal threshold
proportional to $u$ and then sending its proportionality factor to zero
would not justify this result.

\section{Scope and reproducibility}
The result assumes a stationary start, zero drift, independent Poisson complete
resets and a fixed finite horizon on the real line. Bound~\eqref{eq:R} may be
loose at moderate thresholds, but vanishes at fixed parameters. It does not
control joint parameter limits, other renewal laws or nonstationary starts.
It is an exact model calculation, with no experimental inference.

The accompanying standard-library Python diagnostics evaluate the coefficient
and compare last-reset and discounted first-hit representations of $q_a$.
They also compare finite-threshold reflection/renewal calculations at several
reset signs, horizons and rates. Mesh refinement is diagnostic, not a certified
discretization or roundoff bound; the analytic inequalities above establish
the asymptotic and its error. Reproduction instructions and the environment
description accompany the source files.

\paragraph{AI assistance.}
OpenAI Codex was used for literature screening, derivation development,
calculation and diagnostic code, separate-context scientific checking, and
manuscript drafting. These AI checks are not external peer review.
Human author review and approval of the scientific content and this disclosure
remain pending for this local draft.

\paragraph{Data availability.}
No experimental data were used. The calculation code and reproduction
instructions are supplied with the local manuscript package.

\bibliographystyle{plain}
\bibliography{references}
\end{document}
